% Part:sets-functions-relations
% Chapter: size-of-sets
% Section: pairing

\documentclass[../../../include/open-logic-section]{subfiles}

\begin{document}

\olfileid{sfr}{siz}{pai}
\olsection{Fungsi Pasangan lan Kode}

\begin{explain}
Metodhe zig-zag Cantor ndadekake enumerabilitas $\Nat^n$ katon cetha.
Nanging saiki gatekna larikan kang nggambarake $\Nat^2$. Kanthi ngetutake
garis zig-zag lan ngetung posisine, kita bisa mriksa
yen $\tuple{1,2}$ digandhengake karo wilangan~$7$. Nanging bakal
luwih becik yen iki bisa dietung kanthi luwih langsung. Tegese, bakal
migunani yen ana \emph{invers} enumerasi zig-zag,
$g\colon \Nat^2 \to \Nat$, saengga
\[
g(\tuple{0,0}) = 0, \;
g(\tuple{0,1}) = 1, \;
g(\tuple{1,0}) = 2, \; \dots,
g(\tuple{1,2}) = 7, \; \dots
\]
Iki bakal ngidini kita ngetung kanthi pas ing ngendi $\tuple{n, m}$ katon
ing enumerasi.

Sejatine, $g$ bisa langsung ditetepake kanthi rong pangamatan. Kapisan:
yen baris kaping $n$ lan kolom kaping $m$ ngemot nilai~$v$, mula
baris kaping $(n+1)$ lan kolom kaping $(m-1)$ ngemot nilai $v + 1$.
Kapindho: baris kapisan enumerasi ngemot wilangan segitiga,
wiwit $0$, $1$, $3$, $6$, lan sateruse. Wilangan segitiga kaping $k$
yaiku gunggung wilangan natural kang $< k$, kang bisa dietung kanthi
$k(k+1)/2$. Kanthi nggabungake rong pangamatan iki, gatekna
fungsi:
\[
  g(n,m) = \frac{(n+m+1)(n+m)}{2} + n
\]
Kita kerep mung nulis $g(n, m)$ tinimbang $g(\tuple{n, m})$ supaya
luwih gampang diwaca. Rumus iki njaluk kita nemtokake wilangan segitiga
kaping $(n+m)^\text{}$ dhisik, banjur nambah $n$.
Rumus iki ngisi larikan persis kaya kang dikarepake. Mligine,
pasangan $\tuple{1, 2}$ dipetakake menyang $\frac{4 \times 3}{2} +
1 = 7$.

Fungsi $g$ iki \emph{invers} enumerasi sawijining himpunan pasangan.
Fungsi kaya mangkono diarani \emph{fungsi pasangan}.
\end{explain}

\begin{defn}[Fungsi pasangan]
  Fungsi $f\colon A \times B \to \Nat$ iku \emph{fungsi pasangan}
  aritmetis yen $f$ injektif. Kita uga nyebut $f$
  \emph{ngodeake} $A \times B$, lan $f(x,y)$ iku
  \emph{kode} kanggo $\tuple{x,y}$.
\end{defn}

\begin{explain}
Fungsi pasangan bisa digunakake kanggo ngodeake, upamane, pasangan wilangan
natural. Kanthi tembung liya, saben \emph{pasangan} anggota bisa diwakili
nganggo \emph{siji} wilangan. Kanthi invers fungsi pasangan,
kita bisa \emph{ngudhari kode} wilangan iku, yaiku nemtokake
pasangan kang diwakili.
\end{explain}

\begin{prob}
Wenehana enumerasi himpunan kabeh wilangan rasional kang ora negatif.
\end{prob}

\begin{prob}
Tuduhna yen $\Rat$ iku !!{enumerable}. Elinga yen saben wilangan rasional
bisa ditulis minangka pecahan $z/m$ kanthi $z \in \Int$, $m \in \Nat^+$.
\end{prob}

\begin{prob}
Tetepna enumerasi $\Bin^*$.
\end{prob}

\begin{prob}
Elinga saka kuliah pambuka logika yen saben tabel kabeneran
ngandharake fungsi kabeneran. Kanthi tembung liya, fungsi kabeneran
yaiku kabeh fungsi saka $\Bin^k \to \Bin$ kanggo sawijining~$k$.
Buktekna yen himpunan kabeh fungsi kabeneran bisa didhaptar.
\end{prob}

\begin{prob}
Tuduhna yen himpunan kabeh subhimpunan winates saka sembarang
himpunan tanpa wates kang !!{enumerable} iku !!{enumerable}.
\end{prob}

\begin{prob}
Subhimpunan $\Nat$ diarani \emph{kowinates} yen lan mung yen dadi
komplemen sawijining himpunan winates $\Nat$; yaiku $A \subseteq \Nat$
kowinates yen lan mung yen $\Nat\setminus A$ winates. Upamakna $I$ himpunan kang
!!{element}sne persis kabeh subhimpunan winates lan kowinates saka $\Nat$.
Tuduhna yen $I$ iku !!{enumerable}.
\end{prob}

\begin{prob}
Tuduhna yen gabungan kang !!{enumerable} saka himpunan-himpunan kang !!{enumerable}
iku !!{enumerable}. Tegese, yen $A_1$, $A_2$, \dots{} iku himpunan,
lan saben $A_i$ iku !!{enumerable}, mula gabungan $\bigcup_{i=1}^\infty
A_i$ saka kabeh himpunan mau uga !!{enumerable}. [Cathetan: iki angel!]
\end{prob}

\begin{prob}
Upamakna $f \colon A \times B \to \Nat$ sembarang fungsi pasangan.
Tuduhna yen invers saka $f$ iku enumerasi $A \times B$.
\end{prob}

\begin{prob}
Temtokna fungsi kang ngodeake $\Nat^3$.
\end{prob}

\end{document}
